\documentclass[withoutpreface,bwprint]{cumcmthesis}
\usepackage{url}
\usepackage{float}
\usepackage{longtable}
\usepackage{multirow}
\usepackage{booktabs}
\usepackage{array}
\usepackage{listings}
\usepackage{xcolor}
\usepackage{amsmath,amssymb,amsfonts}
\newtheorem{proposition}{\heiti 命题}
\newenvironment{pf}{\par\noindent\textbf{证明}\quad}{\hfill$\blacksquare$\par}
\usepackage{graphicx}
\usepackage{subcaption}
\usepackage{algorithm}
\usepackage{algpseudocode}
\usepackage{tikz}
\usetikzlibrary{arrows.meta,positioning,shapes.geometric,fit,backgrounds,calc}
\usepackage{pgfplots}
\pgfplotsset{compat=1.18}

\setlength{\emergencystretch}{2.5em}
\setlength{\textfloatsep}{7pt plus 2pt minus 2pt}
\setlength{\floatsep}{7pt plus 2pt minus 2pt}
\setlength{\intextsep}{7pt plus 2pt minus 2pt}
\setlength{\abovedisplayskip}{5pt plus 2pt minus 2pt}
\setlength{\belowdisplayskip}{5pt plus 2pt minus 2pt}
\usepackage{etoolbox}
\AtBeginEnvironment{algorithmic}{\scriptsize}
\renewcommand*{\baselinestretch}{1.10}
\AtBeginEnvironment{thebibliography}{\scriptsize}
\apptocmd{\thebibliography}{\setlength{\itemsep}{0.5pt}\setlength{\parskip}{0pt}}{}{}

\lstset{
  basicstyle=\ttfamily\scriptsize,
  breaklines=true,
  frame=single,
  framexleftmargin=2pt,
  numbers=left,
  numberstyle=\tiny\color{gray},
  keywordstyle=\color{blue!70!black},
  commentstyle=\color{green!45!black},
  stringstyle=\color{red!60!black},
  showstringspaces=false,
  columns=flexible,
  extendedchars=true,
  tabsize=2,
  captionpos=b
}
\renewcommand{\lstlistingname}{代码}

\title{无线电干扰源的快速自动定位与清除}
\date{}

\begin{document}

\maketitle

\begin{abstract}
本文为搭载便携式测向机的机器狗建立“几何定位—站位优化—在线搜索规划—可证终止”的完整模型体系，并按附件 1、附件 2 复现了本地模拟环境，完成大规模演练与三次正式测试。

\textbf{问题一}：把带 $\pm1^{\circ}$ 误差的示向度解释为以检测点为顶点的角楔，定位区域即各角楔之交。证明其为凸多边形、给出 $O(n^2)$ 顶点枚举算法与基于回收锥的有界性判据，用旋转卡壳法求直径并与一阶解析式 $D\!\approx\!\sqrt{w_1^2\!+\!w_2^2\!+\!2w_1w_2\cos\gamma}/\sin\gamma$ 互证（细长区域偏差 $<0.1\%$）。对“以直径为直径的圆能否覆盖”的回答是\textbf{一般不能}：本文证明任何直径为 $D$ 的覆盖圆必与直径圆重合（提法良定义），给出 Thales 充要判据，并在 6 万组随机算例中测得不能覆盖的比例为 $1.13\%$（2 站）$\sim3.70\%$（5 站），最大超出半径仅 $0.18\%$，同时给出最小包围圆作为正确替代。

\textbf{问题二}：注意到误差“地点固定、重复检测无效”，将选点建模为\textbf{极大极小优化}：在可行源扇形上最小化最坏情形定位区域直径，并加“保证第二次能收到信号”的约束。解得 $b^{*}=1004.1$ m、$\varphi^{*}=36.8^{\circ}$（相对示向度偏转）、$J^{*}=134.0$ m，且与检测点位置无关。候选区域为以 $S_1$ 为心、半径约 1000 m、张角 $\pm40^{\circ}$ 的圆弧带。相对朴素垂直策略，平均定位区域直径由 96.6 m 降到 55.1 m；沿示向度前进则 25.7\% 的算例定位区域无界；4000 次 Monte-Carlo 无一次越界。

\textbf{问题三}：提出“普查—机会式交会—射线归航—可证终止”策略。由\textbf{七点覆盖定理}（中心加正六边形顶点，$\rho=1280$ m）保证发现全部全向源，由“无信号点集 $R_{\min}$-覆盖目标域”给出可证明的终止判据；位置估计取问题一定位区域的最小包围圆，清除采用步长几何收缩、过冲减半的射线归航与末段六边形清除图案。60 组演练完成率 \textbf{100\%}，平均定位清除时间 \textbf{326.8} s；三次正式测试清除 15/15、12/12、10/10，平均时间 257.5 s、392.0 s、389.1 s。

\textbf{问题四}：定向源造成“\textbf{半平面黑洞}”，单点无信号无法排除干扰源。本文给出\textbf{认证判据}（候选位置落在其 $R_{\min}$ 邻域内无信号点凸包内部时任何定向方向都无法隐藏）与\textbf{射线逼近引理}（沿曾收到信号的射线逼近，全程不丢信号），并把站位加密为间距 1100 m 的三角格点。\textbf{末段定位}采用垂向截断测量（补一次垂直于视线的观测，使最小包围圆半径由 77 m 降到 1 m）与沿射线自适应站位（保证贴边定向源仍在受光侧）。60 组混合演练完成率 \textbf{100\%}，平均定位清除时间 \textbf{597.0} s；三次正式测试清除 14/14、12/12、14/14，平均时间 503.5 s、714.4 s、519.5 s（均值 579.2 s）。

论文另给出多方法互证、残差与不确定度一致性检验、参数灵敏度分析与外部文献对照（最优交会角 $90^{\circ}$、覆盖密度 $2\pi/3\sqrt3$ 等），并讨论模型优缺点与推广。

\keywords{\textbf{交会定位} \quad \textbf{极大极小站位优化} \quad \textbf{圆覆盖} \quad \textbf{在线路径规划} \quad \textbf{可证终止判据}}
\end{abstract}

\thispagestyle{empty}

\section{问题重述}

\subsection{问题背景}

无线电干扰源的定位与清除是无线电频谱管理的核心任务。《中华人民共和国无线电管理条例》（国令第 672 号）第五十七条明确规定无线电监测站承担“对无线电信号实施监测，查找无线电干扰源和未经许可设置、使用的无线电台（站）”的职责\cite{reg672}。工程上干扰源定位分为初步监测、区域性排查与精准定位三阶段，前两阶段可由固定站、监测车或无人机完成，而受遮挡与低功率源制约，最终精确定位仍需技术人员携便携测向机徒步排查，效率低、风险高。

本题要求为某公司拟研发的“搭载于机器狗的干扰源自动定位清除系统”建立数学模型。系统工作环境为：半径 1800 m 的圆形目标区域，圆心为坐标原点，$x$ 轴指向正东、$y$ 轴指向正北，单位为米；区域内干扰源个数未知。设备与规则见附录 1 与附录 2，要点如下：

\begin{itemize}
  \item 每个干扰源的频道互不相同，取自 20 个频道 $\{1,\dots,20\}$，不随时间变化；不同干扰源信号互不影响。
  \item 干扰源分为\textbf{全向}（覆盖 $360^{\circ}$）与\textbf{定向}（覆盖以定向方向为中心的 $180^{\circ}$ 扇形，含边界）两类；在覆盖角内信号均匀辐射，场强仅随距离衰减。
  \item 测向机在某点测得的\textbf{示向度}为接收场强最大方向的方位角，范围 $[0^{\circ},360^{\circ})$；由于本地电磁环境影响，示向度与真实方位角存在误差，误差在 $[-1^{\circ},1^{\circ}]$ 内；\textbf{同一地点电磁环境固定，重复检测不改变误差}，只有不同地点才呈现统计规律。
  \item 每个干扰源存在最大可被有效接收距离（有效接收半径），取值在 $1000\sim1500$ m 之间且各不相同。
  \item 交会定位法：两个检测点的两条示向度方向线及其两侧各偏 $1^{\circ}$ 的虚线围成的四边形称为\textbf{定位区域}。
  \item \textbf{时间与清除}：频道切换 1 s；一次检测（停下、切换、调天线、读稳定示向度）5 s；机器狗沿直线以 5 m/s 移动。光学探测仪仅在距干扰源 $\le20$ m 时可精确定位并启动激光枪清除（定位 3 s、清除 2 s）；20 m 内无该频道源时只耗 3 s（返回“未发现”）。检测点距源 $\le5$ m 且在其覆盖角内时信号过强而无法给出示向度（返回“near”），可直接清除。
\end{itemize}

\subsection{需要解决的问题}

\begin{enumerate}
  \item \textbf{问题一}：已知若干检测点坐标及某干扰源关于这些检测点的示向度，给出采用交会定位法计算\textbf{多边形定位区域直径}（区域内任意两点距离的最大值）的算法；并回答“以定位区域直径为直径的圆能否覆盖此定位区域”。
  \item \textbf{问题二}：已知某全向干扰源在一个检测点处测得的示向度，给出\textbf{第二个检测点的选择策略}以获得较好的定位效果，并给出第二检测点的\textbf{候选区域}。
  \item \textbf{问题三}：目标区域内全为全向干扰源，总数在 $10\sim16$ 之间但具体个数未知。机器狗从原点出发、测向机初始频道为 1，需在保证全部清除的前提下使完成任务时间尽量短：制定自动搜索定位及清除策略并设计算法，通过模拟器演练测试检验，统计\textbf{被清除干扰源个数的比例}与\textbf{平均定位清除时间}（定位清除总时间含移动、频道切换、检测、光学精确定位与清除耗时），随后进行三次正式测试并按表 \ref{tab:formal} 格式给出结果，导出日志文件。
  \item \textbf{问题四}：区域内既有全向又有定向干扰源，总数仍为 $10\sim16$，定向源的个数与定向方向均未知，其余条件同问题三，给出检测与清除的策略与算法并演练测试、三次正式测试。
\end{enumerate}

\begin{table}[H]
\centering
\caption{问题 3/4 正式测试结果表格式（题目表 1）}\label{tab:formal}
\begin{tabular}{cccc}
\toprule
测试案例编码 & 清除干扰源个数 & 平均定位清除时间 & 程序运行时间\\
\midrule
测试 1 &&& \\
测试 2 &&& \\
测试 3 &&& \\
\bottomrule
\end{tabular}
\end{table}

\section{问题分析}

\subsection{总体分析框架}

四个问题在逻辑上层层递进，构成一条完整的技术链（图 \ref{fig:frame}）：

\begin{center}
\resizebox{\textwidth}{!}{%
\begin{tikzpicture}[
  font=\small,
  box/.style={draw, rounded corners=3pt, align=center, inner sep=6pt,
              minimum height=1.50cm, minimum width=2.58cm, fill=blue!6,
              line width=0.7pt},
  core/.style={draw, rounded corners=3pt, align=center, inner sep=6pt,
               fill=orange!14, line width=0.7pt},
  lbl/.style={font=\scriptsize, text=black!70, inner sep=1.5pt},
  ar/.style={-{Stealth[length=2.2mm]}, line width=0.8pt},
  dar/.style={-{Stealth[length=2.2mm]}, line width=0.8pt, dashed}
]
% --- four problem boxes -------------------------------------------------
\node[box] (q1) {\textbf{问题一}\\[1pt] 定位区域 $\mathcal{R}$\\[1pt] 直径 $D$、圆覆盖判据};
\node[box, right=1.50cm of q1] (q2) {\textbf{问题二}\\[1pt] 第二检测点\\[1pt] minimax 站位优化};
\node[box, right=1.50cm of q2] (q3) {\textbf{问题三}\\[1pt] 全向源在线\\[1pt] 搜索—定位—清除};
\node[box, right=1.50cm of q3] (q4) {\textbf{问题四}\\[1pt] 含定向源\\[1pt] 检测与认证策略};

% --- chain arrows with labels placed ABOVE the arrows -------------------
\draw[ar] (q1) -- node[lbl, above=2pt] {不确定度度量} (q2);
\draw[ar] (q2) -- node[lbl, above=2pt] {站位规则} (q3);
\draw[ar] (q3) -- node[lbl, above=2pt] {覆盖与认证} (q4);

% --- shared kernel ------------------------------------------------------
\node[core, minimum width=14.82cm] (core)
     at ($(q1.south)!0.5!(q4.south) + (0,-1.15cm)$)
     {\textbf{统一内核}：角度楔交会 $\Rightarrow$ 定位区域（凸多边形）$\Rightarrow$ 最小包围圆半径作不确定度 $\sigma$};

% --- kernel feeds every problem ----------------------------------------
\foreach \q in {q1, q2, q3, q4}{%
  \draw[dar] (core.north -| \q.south) -- (\q.south);
}

\end{tikzpicture}}
\captionof{figure}{四个问题的逻辑关系：问题一给出几何内核与不确定度度量，问题二给出站位规则，
问题三、四是在线决策问题；统一内核（下方橙色框）为四个问题共用}
\label{fig:frame}
\end{center}

问题一、二是\textbf{几何与最优设计问题}：给定带界误差的方位测量，刻画所有与测量相容的干扰源位置集合（定位区域），并以此为核心指标优化观测几何。问题三、四则是\textbf{在线决策与规划问题}：干扰源的位置、频道、个数均需在行进中逐步发现，机器人必须在“继续搜索”与“前往清除已知目标”之间分配时间，并且必须在有限时间内\textbf{自行判断任务是否完成}。三个层次的关键联系在于：问题一给出的定位区域及其最小包围圆，恰好是问题三、四中位置估计的\textbf{不确定度}；问题二的站位规则给出了获取第二个方位的最优几何；而问题三、四的终止判据则依赖于“探测圆覆盖”这一纯几何结论。

\subsection{问题一分析}

关键在于把“示向度 $\pm1^{\circ}$”正确地翻译成几何约束。设检测点 $S_i$ 处测得的示向度为 $\hat\theta_i$，则干扰源真实方位角 $\theta_i$ 满足 $|\theta_i-\hat\theta_i|\le 1^{\circ}$，等价于干扰源位于以 $S_i$ 为顶点、角平分线方向为 $\hat\theta_i$、半张角 $1^{\circ}$ 的\textbf{角楔} $W_i$ 内。于是“与全部测量相容的位置集合”即
\[
\mathcal{R}=\bigcap_{i=1}^{n} W_i,
\]
这正是附录 2 图 2 中红色四边形的推广（$n=2$ 时退化为四边形）。由此可得三条性质：$\mathcal{R}$ 是凸的（凸集之交）；每个 $W_i$ 是无界凸锥，故 $\mathcal{R}$ 可能\textbf{有界、无界或为空}；$\mathcal{R}$ 的边界必由若干条示向度边界直线支撑，因此其顶点必为两条边界直线的交点。这三条性质直接给出 $O(n^2)$ 的顶点枚举算法与基于回收锥的有界性判据。

第二个问题“以定位区域直径为直径的圆能否覆盖该区域”容易被直觉误导。由 Thales 定理，点 $P$ 落在以 $AB$ 为直径的闭圆盘内当且仅当 $\angle APB\ge90^{\circ}$，即 $(P-A)\cdot(P-B)\le0$；由于 $\mathcal{R}$ 是凸多边形，只需对顶点逐一检验即可判定。进一步，若存在某个直径为 $D$ 的圆覆盖 $\mathcal{R}$，则该圆必同时包含直径端点 $A,B$，而同时包含 $A,B$ 且半径为 $D/2$ 的圆唯一，即以 $AB$ 为直径的圆；因此“以直径为直径的圆能否覆盖”与“是否存在直径为 $D$ 的覆盖圆”是同一问题，提法良定义。对一般凸集答案是\textbf{否}（等边三角形、细长四边形都是反例），对本文的定位区域需要定量回答，故第三、四节用解析论证与 6 万次随机算例给出结论及其量级。

\subsection{问题二分析}

题面只给“一个检测点处测得的示向度”，若把第二检测点选在任意位置，定位区域可能很大甚至无界。选择策略必须只依赖于 $\hat\theta_1$、可行源集合与约束条件。逐层分析：

\textbf{(1) 信息来源的全部刻画}。信号在 $S_1$ 被收到 $\Rightarrow$ 距离不超过有效接收半径 $\le1500$ m；收到的是 direction 而非 near $\Rightarrow$ 距离 $>5$ m；再加 $G$ 在目标域内与方位角 $\pm1^{\circ}$ 约束，可行源集合 $\mathcal{G}$ 是一段长 1500 m、宽 $\pm1^{\circ}$ 的扇形（图 \ref{fig:q2geo}）。

\textbf{(2) 评价指标}。以第二检测点 $S_2$ 观测后所得定位区域的直径 $D$ 作为定位精度的度量（直径是区域尺寸的极大极小刻画，且与问题一直接衔接）。

\textbf{(3) 设计准则}。由于 $G$ 在 $\mathcal{G}$ 中未知，取\textbf{极大极小（minimax）准则}：
\begin{equation}
S_2^{*}=\arg\min_{S_2}\ \max_{G\in\mathcal{G}}\ \max_{|\delta_2|\le1^{\circ}}\ D\bigl(S_1,S_2,G,\delta_2\bigr).
\label{eq:q2minimax}
\end{equation}
同时必须保证第二次检测确实能收到信号，否则策略失效，故施加约束：对一切 $G\in\mathcal{G}$ 有 $|S_2-G|\le R_{\min}=1000$ m（取最坏接收半径，故为\textbf{保证}探测）。

\textbf{(4) 求解思路}。由一阶解析式 $D\propto\sqrt{w_1^2+w_2^2+2w_1w_2\cos\gamma}/\sin\gamma$（$w_i=2\varepsilon d_i$ 为两侧向线在源处的横向不确定度，$\gamma$ 为交会角）可见，$D$ 由两个因素控制：交会角 $\gamma$ 应接近 $90^{\circ}$，且到源的距离应尽量小。前者要求 $S_2$ 落在以 $S_1\widehat{G}$ 为直径的圆附近，后者又要求靠近源，二者与“必须落在 $\mathcal{G}$ 的 $1000$ m 邻域内”这一可行性约束相互制约，因此最优解存在且不退化。

\textbf{(5) 为什么不能沿示向度方向前进}。若第二检测点沿示向度方向前进，两条方位线近似平行，回收锥非空，定位区域无界；本文用 3000 次 Monte-Carlo 定量说明这一点（第 6.8 节）：沿示向度前进 800 m 时 25.7\% 的算例定位区域无界。

\subsection{问题三分析}

\textbf{(1) 与问题一、二的差别}。问题一、二是“测量已知、位置待定”，问题三是“连有哪些干扰源都不知道”，本质上是\textbf{在线（online）搜索与路径规划}问题：频道 $1\sim20$ 中只有 $10\sim16$ 个被占用，且具体个数未知，机器人必须通过检测逐一发现。

\textbf{(2) 时间结构}。虚拟时间由四部分构成：移动（$1$ s 对应 $5$ m）、频道切换（两频道之间恒为 1 s）、检测动作（恒 5 s）、清除动作（命中 5 s、未命中 3 s）。程序运行时间上限 20 min 只在真实时间维度上约束（本地模拟中每次请求的往返时间约 1 ms，故真实时间不构成瓶颈），因此优化目标是\textbf{虚拟时间}。由量纲分析可知：移动 500 m 的时间（100 s）相当于 20 次检测（每次 5 s），因此\textbf{减少无效移动比减少检测次数更重要}，但盲目减少检测又会导致漏检与无法终止，二者需要统筹。

\begin{figure}[H]
\centering
\includegraphics[width=0.74\textwidth]{../figures/fig_simulator_timing.png}
\caption{四条指令中各动作的虚拟时间消耗（按附件 1 的算例：移动 500 m 用 100 s，检测 5 s，切换频道 1 s，清除未命中 3 s、命中 5 s）}
\label{fig:timing}
\end{figure}

\textbf{(3) 三个必须在模型中解决的子问题}。
\begin{itemize}
  \item \emph{保证发现}：如何安排检测位置，使得无论干扰源在区域内何处、有效接收半径取何值（$\ge1000$ m），都至少被检测到一次？这是纯几何的\textbf{圆覆盖问题}：用半径 $R_{\min}=1000$ m 的圆覆盖半径 1800 m 的圆域所需最少圆数。
  \item \emph{可证终止}：干扰源个数未知，机器人必须能\textbf{证明}某些频道不存在干扰源，否则永远无法判断“是否已清除全部”。
  \item \emph{在线路径}：已知目标需清除、未知区域需搜索，如何在二者之间分配时间（在线 TSP / 最小延迟问题的工程化处理）。
\end{itemize}

\textbf{(4) 终止判据的思路}。若某频道在若干位置处均返回“no\_signal”，且这些位置到目标域内任意候选位置的距离都不超过 $R_{\min}$，则该频道必然不存在未清除的干扰源（否则该处一定能收到信号）。因此“无信号点集是否 $R_{\min}$-覆盖目标域”给出了一个\textbf{可证明}的终止条件，且它与“保证发现”用的是同一套覆盖几何，可实现“搜索即认证”。

\subsection{问题四分析}

定向源把问题从“覆盖”提升到“可辨识”层次。设定向源位于 $q$、定向方向为 $u$，则检测点 $p$ 能收到信号的条件是 $|p-q|\le R_{\rm eff}$ \textbf{且} $u\cdot(p-q)\ge0$。于是：

\textbf{(1) 半平面黑洞}。固定 $p$，对任意 $q$，总可以选取 $u$ 使 $u\cdot(p-q)<0$，即让该点落在覆盖角之外。因此\textbf{单点无信号永远不能排除定向源}，问题三的“覆盖即认证”判据在问题四中失效，必须强化。

\textbf{(2) 强化后的认证判据}。若在 $q$ 的 $R_{\min}$ 邻域内已有一组无信号检测点 $\{p_j\}$，则“存在某个定向方向躲过全部检测”等价于“存在半平面包含全部 $p_j-q$”。由凸集分离定理，后者不成立当且仅当\textbf{原点位于 $\{p_j-q\}$ 凸包内部}。这给出一个可在常数时间内检验的判据，也解释了为什么问题四的搜索站位必须比问题三密（需要从多个方向包围候选位置）。

\textbf{(3) 定向源的可逼近性}。定向源的另一困难是接近过程中可能“丢失信号”。本文证明：若在 $p$ 处收到过信号，则线段 $[p,q]$ 上任意点 $p'$ 都满足 $u\cdot(p'-q)\ge0$（因为 $p'-q=\lambda(p-q),\ 0\le\lambda\le1$），即\textbf{沿曾收到信号的射线走向源，信号不会丢失}。这一引理使问题四的末段逼近可以复用问题三的射线归航算法，只需在丢失信号时回退到上一次听到信号的位置。

\textbf{(4) 区域外采样的必要性}。贴边且朝外辐射的定向源，其覆盖角完全落在目标域之外，域内任何检测点都收不到信号；此时必须（题目允许机器狗提交区域外的坐标）到区域外补测才能完成认证。本文将其作为可选机制并定量分析其时间代价。

\section{模型假设}

\begin{enumerate}
  \item \textbf{平面假设}：由附录 2(7)，测向机与全部干扰源位于同一平面，几何问题均为二维问题。
  \item \textbf{误差的有界性与地点确定性}：示向度误差 $\delta\in[-1^{\circ},1^{\circ}]$；误差是检测点位置的确定性函数（同一地点重复检测结果完全相同），不同地点的误差相互独立且可在区间内任意取值。因此多次重复检测\textbf{不提供}额外信息，只有更换检测点才提高精度。
  \item \textbf{接收条件}：频道 $c$ 的干扰源在检测点 $p$ 被收到，当且仅当 $|p-q_c|\le R_c$（全向源）或 $|p-q_c|\le R_c$ 且 $u_c\cdot(p-q_c)\ge0$（定向源），其中 $R_c\in[1000,1500]$ m。反向结论“收不到”只说明上述条件被破坏，三种原因（不存在、超距、不在覆盖角内）无法区分，模型必须显式处理这种多义性。
  \item \textbf{动作与计时}：移动耗时 $=$ 直线距离 $/5$；任意两频道切换耗时 1 s；一次检测固定 5 s；清除命中 5 s、未命中 3 s；$\text{/clear}$ 不改变测向机当前频道、不产生频道切换耗时。
  \item \textbf{清除条件}：$p$ 处对频道 $c$ 的清除成功当且仅当 $|p-q_c|\le20$ m，\textbf{与干扰源类型和覆盖角无关}；同一干扰源只能被清除一次。
  \item \textbf{统计独立}：各干扰源的位置、频道、类型、有效接收半径相互独立；位置在目标圆域内均匀分布（用于演练与灵敏度分析）。
  \item \textbf{通信可靠}：与模拟器的通信满足附录 2 的协议；网络短暂异常时按“同 request\_id 重试”处理，不影响虚拟时间记账。
\end{enumerate}

\section{符号说明}

\begin{center}
\begin{longtable}{cl}
\toprule
符号 & 含义\\
\midrule
\endhead
$S_i=(x_i,y_i)$ & 第 $i$ 个检测点（站位）坐标，单位 m\\
$\hat\theta_i,\ \theta_i$ & $S_i$ 处测得的示向度与干扰源的真实方位角，单位度\\
$\varepsilon$ & 示向度误差上界，本题 $\varepsilon=1^{\circ}$\\
$\delta_i$ & $S_i$ 处的示向度误差，$|\delta_i|\le\varepsilon$\\
$W_i$ & 由 $S_i$ 的示向度确定的角楔（$\pm\varepsilon$ 锥）\\
$\mathcal{R}$ & 定位区域，$\mathcal{R}=\bigcap_i W_i$\\
$D$ & 定位区域直径，$D=\max_{P,Q\in\mathcal{R}}|P-Q|$\\
$R_{\mathrm{MEC}}$ & 定位区域最小包围圆（MEC）半径\\
$\gamma$ & 交会角，即干扰源处两条方位视线的夹角\\
$w_i=2\varepsilon d_i$ & 第 $i$ 条方位线在源处的横向不确定度（弧度为 $\varepsilon$）\\
$\mathcal{G}$ & 与已有测量相容的可行干扰源集合\\
$R_{\rm eff},R_c$ & 干扰源的有效接收半径；$R_{\min}=1000$ m，$R_{\max}=1500$ m\\
$q_c,u_c$ & 频道 $c$ 的干扰源位置与定向方向（全向源无 $u_c$）\\
$J(S_2)$ & 站位的极大极小指标：$\max_{G\in\mathcal{G}}D$\\
$b,\varphi$ & 第二检测点相对 $S_1$ 的距离与相对示向度方向的偏转角\\
$\sigma$ & 位置估计的不确定度半径（取定位区域 MEC 半径）\\
$T,\bar T$ & 定位清除总时间、平均定位清除时间\\
$N_c$ & 频道 $c$ 被清除的状态（未清除/已清除/已认证不存在）\\
\bottomrule
\end{longtable}
\end{center}
